% p.64 — الباب التاسع والعشرون في الاصطلاحات والتركيبات (تتمة)
\mspage{64}
% Layout: one boxed table of four columns (two pairs, right pair first).  The upper part
% of the right pair continues the terminology list (43–56, plain digits).  Below a rule,
% the تركيب examples begin: each example sentence carries a circled number (here (n)),
% and under it each word is analysed; the narrow right-hand column of each pair repeats
% the grammatical label alone, the wide column gives «word + label».  The analysis of
% example (2) runs from the foot of the right pair to the head of the left pair.
% Printed page number ۶۴ at top left.  Layout: the two pairs are set side by side (right
% pair first) in one `tables` row, as in the lithograph; the tables are split exactly at
% the printed column break (right pair: 43–56, rule, analysis (1)–(2) down to
% «في الدار متعلق»; left pair: «موجود متعلق محذوف» onwards).
\begin{tables}
\begin{sarf}{|r|r|}
(43) فعل مقاربة & كَادَ زَيْدٌ أَنْ يَخْرُجَ\\
(44) فعل مدح & نِعْمَ زَيْدٌ\\
(45) فعل ذم & بِئْسَ زَيْدٌ\\
(46) فعل يقين & عَلِمَ زَيْدٌ عَمْرًا عَالِمًا\\
(47) فعل شك & ظَنَّ زَيْدٌ عَمْرًا كَاتِبًا\\
(48) اسم فعل & هَا زَيْدًا\\
(49) اسم استفهام & مَتٰى تَذْهَبُ\\
(50) اسم شرط & مَتٰى تَذْهَبْ أَذْهَبْ\\
(51) جملة اسمية & زَيْدٌ ضَارِبٌ . زَيْدٌ يَضْرِبُ\\
(52) جملة فعلية & ضَرَبَ زَيْدٌ . اُنْصُرْ زَيْدًا\\
(53) جملة خبرية & زَيْدٌ ضَارِبٌ ضَرَبَ زَيْدٌ\\
(54) جملة انشائية & قُمْ لَا تَقُمْ هَلْ تَقُوْمُ\\
(55) جملة شرطية & اِنْ ضَرَبْتَ زَيْدًا . فَأَضْرِبُكَ\\
(56) جملة جزائية & اِنْ ضَرَبْتَ زَيْدًا فَأَضْرِبُكَ\\\hline
% Analysis section (below a rule, still in the right pair): column 1 = label (as printed
% in the narrow column), column 2 = text.
 & (1) مَرَرْتُ بِزَيْدٍ\\
حرف جر & باء حرف جر\\
مجرور & زيد مجرور\\
متعلِّق & بزيد متعلق\\
مُتعلَّق & مررت متعلق\\
 & (2) زَيْدٌ فِي الدَّارِ زَيْدٌ مَوْجُوْدٌ فِي الدَّارِ\\
مُتعلِّق & في الدار متعلق\\
\end{sarf}
% ---- left pair (continues the analysis of example (2)) ----
\begin{sarf}{|r|r|}
متعلَّق محذوف & موجود متعلق محذوف\\
 & (3) رَأَيْتُ غُلَامَ زَيْدٍ يَوْمَ الْجُمُعَةِ\\
مضاف & غلام مضاف\\
مضاف اليه & زيد مضاف اليه\\
ظرف مضاف & يوم ظرف مضاف\\
مضاف اليه & الجمعة مضاف اليه\\
متعلِّق & يوم متعلِّق\\
مُتعلَّق & رَأَيْتُ متعلَّق\\
 & (4) هٰذَا رَجُلٌ عَالِمٌ\\
موصوف & رجل موصوف\\
صفة & عالم صفة\\
مبتدى & هذا مبتدى\\
خبر & رجل خبر\\
 & (5) زَيْدٌ قَائِمٌ\\
مبتدى & زيد مبتدى\\
خبر & قائم خبر\\
 & (6) زَيْدٌ ضَرَبَ عَمْرًا\\
مبتدى & زيد مبتدى\\
جملة خبر & ضرب عمرا جملة خبر\\
 & (7) جَاءَ رَجُلٌ ضَرَبَ غُلَامُهُ عَمْرًا\\
موصوف & رجل موصوف\\
\end{sarf}
\end{tables}
